\documentclass[11pt]{article}

\usepackage[margin=1in]{geometry}
\usepackage{amsmath,amssymb,amsthm,mathtools}
\usepackage{enumitem}
\usepackage{xcolor}
\usepackage{hyperref}

\title{exp0526: Pure 0522-Style Closed-Loop Controller Learning\\
for Age-Structured Predator--Prey PDEs}
\author{exp0526 notes}
\date{\today}

\newtheorem{problem}{Problem}
\newtheorem{objective}{Objective}
\newtheorem{remark}{Remark}

\newcommand{\R}{\mathbb{R}}
\newcommand{\dd}{\,\mathrm{d}}
\newcommand{\inner}[2]{\left\langle #1,#2\right\rangle}
\newcommand{\softplus}{\operatorname{softplus}}
\newcommand{\bluequestion}[1]{\textcolor{blue}{\textbf{Question.} #1}}

\begin{document}

\maketitle

\section{Purpose}

This note reformulates the age-structured predator--prey control setting
from \emph{Lotka--Sharpe Neural Operators for Control of Population PDEs}
as a pure exp0522-style closed-loop controller learning problem
\cite{krstic2026lotkasharpe}.

The first exp0526 version should not learn a Lotka--Sharpe neural
operator and should not approximate the paper's analytic controller.
Instead, the model directly outputs a scalar control signal \(u(t)\),
applies it to the population dynamics, and receives loss from the
resulting distance to the equilibrium state.

The controller interface is
\[
  \big(\eta_1(t),\eta_2(t),1;\ \text{internal self-talk}\big)
  \longmapsto
  u(t).
\]
The ordinary external observation is \((\eta_1,\eta_2)\) together with a
constant pulse \(1\). The self-talk state is internal to the exp0522
model: the previous language output is fed into the next step.

\section{Problem From the Miroslav Paper}

The environment is the two-species age-structured predator--prey PDE
used in the Lotka--Sharpe neural-operator paper and inherited from the
age-structured predator--prey stabilization literature
\cite{krstic2026lotkasharpe,veil2026stabilization}. It is a
McKendrick--von Foerster type age-structured model
\cite{mckendrick1925applications}:
\begin{subequations}
\label{eq:full-pde}
\begin{align}
  \partial_t x_1(a,t)+\partial_a x_1(a,t)
  &=
  -x_1(a,t)
  \left[
    \mu_1(a,t)+u(t)
    +\int_0^A g_1(\alpha)x_2(\alpha,t)\dd\alpha
  \right],
  \\
  \partial_t x_2(a,t)+\partial_a x_2(a,t)
  &=
  -x_2(a,t)
  \left[
    \mu_2(a,t)+u(t)
    +\frac{1}{\int_0^A g_2(\alpha)x_1(\alpha,t)\dd\alpha}
  \right],
  \\
  x_i(0,t)
  &=
  \int_0^A k_i(a,t)x_i(a,t)\dd a,
  \qquad i=1,2.
\end{align}
\end{subequations}
Here \(x_i(a,t)>0\) are the age profiles, \(k_i\) are fertility profiles,
\(\mu_i\) are mortality profiles, \(g_i\) are interaction kernels, and
\(u(t)\ge 0\) is the scalar dilution or harvesting input.

The control target is a positive equilibrium profile
\[
  x_i(a,t)\to x_i^*(a,t),
  \qquad i=1,2,
\]
or, equivalently for the reduced controller coordinates, the equilibrium
\[
  \eta_1(t)\to 0,
  \qquad
  \eta_2(t)\to 0.
\]

\section{Equilibrium Error Coordinates}

For each species, the Lotka--Sharpe scalar is defined implicitly by the
classical Lotka--Sharpe age-distribution condition
\cite{sharpe1911problem}:
\begin{equation}
  \int_0^A
  k_i(a,t)
  \exp\left(
    -\int_0^a \big(\mu_i(s,t)+\zeta_i(t)\big)\dd s
  \right)
  \dd a
  =
  1.
  \label{eq:lotka-sharpe}
\end{equation}
The Lotka--Sharpe neural-operator paper uses this quantity to construct
transformed coordinates \cite{krstic2026lotkasharpe}
\[
  \eta(t)=(\eta_1(t),\eta_2(t)).
\]
In the paper's notation,
\begin{equation}
  \eta_i(t)
  =
  \log\left(
    \frac{\inner{\pi_{0,i}(\cdot,t)}{x_i(\cdot,t)}}{x_i^*(0,t)\kappa_i(t)}
  \right).
  \label{eq:eta}
\end{equation}

The exp0526 controller observes \(\eta_1,\eta_2\), not the full hidden
ecological parameter state. The full PDE remains the environment, while
\(\eta\) is the low-dimensional control error exposed to the agent.

\section{What We Change}

The Lotka--Sharpe neural-operator paper keeps the analytic feedback law
and uses learning only
to approximate the implicit Lotka--Sharpe map
\[
  (k,\mu)\mapsto\zeta.
\]
That is not the first exp0526 task.

By contrast, the first exp0526 task is direct closed-loop controller
learning:
\[
  \big(\eta_1(t),\eta_2(t),1;\ m_{t-1}\big)
  \longmapsto
  u(t),
\]
where \(m_{t-1}\) is the model's own previous language output. The model
is judged by what its control does to the environment, not by whether it
matches an analytic control formula.

\begin{problem}[Pure exp0526 closed-loop controller]
Train an exp0522-style self-talk MLP inside the full two-species
age-structured PDE environment \eqref{eq:full-pde}. At each time step the
agent observes \(\eta_1(t)\), \(\eta_2(t)\), and a constant pulse \(1\),
produces a scalar control \(u(t)\), and is optimized through the resulting
future equilibrium error.
\end{problem}

\section{0522 Mechanism Reused}

The exp0522 mechanism is not an external memory module. It is a recurrent
self-talk loop created by feeding the previous language output into the
next MLP step.

At step \(t\), the hidden network computes
\[
  h_t = F_\theta\big(\eta_1(t),\eta_2(t),1,e_{t-1},m_{t-1}\big),
\]
then emits both a control preactivation and a new language state:
\begin{align}
  \tilde u_t &= W_u h_t+b_u,\\
  m_t &= R h_t.
\end{align}
Here \(R\) is the fixed sparse signed readout used by exp0522. The next
step receives \(m_t\). This is the sense in which the model has memory:
the language output from the previous step becomes part of the next step's
input.

The ordinary problem statement remains
\[
  (\eta_1,\eta_2,1)\mapsto u,
\]
because \(m_t\) is an internal model state rather than an externally
measured ecological variable.

\subsection*{Error Channel}

In exp0522 prediction, the error channel is a scalar prediction error.
In exp0526 control, the first version should reinterpret the error channel
as a control-error view:
\[
  e_t = \eta_t = (\eta_1(t),\eta_2(t)).
\]
This is not a control-label error. It is the current distance direction
from equilibrium, made available through the same error-feedback interface
so that the 0522 ablations remain meaningful.

The no-error comparison forces this channel to zero. The no-language
comparison uses \(m_t\equiv 0\). Blink and stutter tests temporarily
remove the error or language pathway during rollout.

\section{Closed-Loop Learning Objective}

The network produces an unconstrained scalar \(\tilde u_t\). The actual
control applied to the environment is
\[
  u_t=\softplus(\tilde u_t),
  \label{eq:softplus-control}
\]
which enforces \(u_t\ge 0\) while preserving smooth gradients.

For one differentiable rollout window of length \(T\), the main objective
is distance to equilibrium:
\begin{equation}
  \mathcal L_\eta
  =
  \frac{1}{T}
  \sum_{t=1}^{T}
  \left(
    \eta_1(t)^2+\eta_2(t)^2
  \right).
  \label{eq:eta-loss}
\end{equation}
A small control-energy regularizer can be added:
\begin{equation}
  \mathcal L
  =
  \mathcal L_\eta
  +
  \lambda_u
  \frac{1}{T}
  \sum_{t=1}^{T}
  u_t^2,
  \qquad
  0<\lambda_u\ll 1.
  \label{eq:total-loss}
\end{equation}

The required gradient path is
\[
  \mathcal L
  \to
  \eta_{t+1}
  \to
  x(\cdot,t+1)
  \to
  u_t
  \to
  \tilde u_t
  \to
  h_t
  \to
  \theta.
\]
Because \(m_t=Rh_t\) is fed to later steps, gradients inside a training
window also flow through the self-talk loop. At window boundaries, the
message and PDE state are carried forward but detached, matching the
continuous-window protocol from exp0522.

\subsection*{Backpropagation Through the Discrete Rollout}

\bluequestion{How exactly does the equilibrium loss produce a gradient
with respect to the controller parameters?}

In implementation, the continuous PDE is replaced by a differentiable
time-step map. For a fixed hidden ecological profile at step \(n\), write
the age-grid state as \(X_n\in\R^{2\times N_a}\) and the controller state
as \(m_{n-1}\). The model and environment form the closed-loop recursion
\begin{align}
  o_n
  &=
  \big[\eta(X_n;\xi_n),\,1\big],
  \\
  (\tilde u_n,m_n)
  &=
  C_\theta(o_n,e_n,m_{n-1}),
  \\
  u_n
  &=
  \softplus(\tilde u_n),
  \\
  X_{n+1}
  &=
  \Phi_{\Delta t}(X_n,u_n;\xi_n),
  \label{eq:closed-loop-discrete-map}
\end{align}
where \(\xi_n=(k_n,\mu_n,g_n,\zeta_n,\pi_{0,n},x_n^*)\) denotes the
environment quantities used at that step. The map
\(\Phi_{\Delta t}\) is the upwind age-grid update together with the
renewal boundary condition. The training loss over one window
\(n=0,\ldots,T-1\) is
\begin{equation}
  \mathcal L_T(\theta)
  =
  \frac{1}{T}
  \sum_{n=1}^{T}
  \|\eta(X_n;\xi_n)\|_2^2
  +
  \lambda_u
  \frac{1}{T}
  \sum_{n=0}^{T-1}
  u_n^2 .
  \label{eq:window-loss}
\end{equation}
Thus the controller is not trained from a target control label. It is
trained by differentiating the consequence of its own control actions
through the rollout.

\bluequestion{What chain rule is being used inside one training window?}

Let \(z_n=(X_n,m_{n-1})\) be the recurrent closed-loop state. Equations
\eqref{eq:closed-loop-discrete-map} define a differentiable recurrence
\[
  z_{n+1}=G_\theta(z_n;\xi_n).
\]
Backpropagation through time computes
\begin{equation}
  \nabla_\theta \mathcal L_T
  =
  \sum_{n=0}^{T-1}
  \left(
    \frac{\partial \ell_n}{\partial u_n}
    \frac{\partial u_n}{\partial \tilde u_n}
    \frac{\partial \tilde u_n}{\partial \theta}
  \right)
  +
  \sum_{n=1}^{T}
  \frac{\partial \ell_n}{\partial X_n}
  \frac{\partial X_n}{\partial \theta},
  \label{eq:bptt-gradient}
\end{equation}
where \(\ell_n=\|\eta(X_n;\xi_n)\|_2^2/T+\lambda_u u_n^2/T\), and
\(\partial X_n/\partial\theta\) is itself obtained recursively from
\[
  \frac{\partial X_{n+1}}{\partial\theta}
  =
  \frac{\partial \Phi_{\Delta t}}{\partial X_n}
  \frac{\partial X_n}{\partial\theta}
  +
  \frac{\partial \Phi_{\Delta t}}{\partial u_n}
  \frac{\partial u_n}{\partial\tilde u_n}
  \frac{\partial \tilde u_n}{\partial\theta}.
\]
The important term is
\(\partial \Phi_{\Delta t}/\partial u_n\): changing the neural control
changes the next population profile, which changes future \(\eta\), which
changes the loss. Modern autodiff computes the same chain rule directly
from the PyTorch tensor graph; no analytic reference controller is needed.

\bluequestion{Do we differentiate through the Lotka--Sharpe root solve?}

In the first implementation, no. The hidden ecological profile
\(\xi_n\), including \(\zeta_n\), \(\pi_{0,n}\), \(\kappa_n\), and the
instantaneous equilibrium profile \(x_n^*\), is treated as environment
truth at step \(n\). These quantities define the measurement
\(\eta(X_n;\xi_n)\), but their root-solve construction is detached from
the computational graph. Gradients are still taken with respect to the
population state \(X_n\), the control \(u_n\), the controller hidden state,
the self-talk message, and the network parameters \(\theta\).

\bluequestion{What happens at continuous-window boundaries?}

Within a window, the graph contains all steps
\[
  (X_0,m_{-1})
  \to
  (u_0,X_1,m_0)
  \to
  \cdots
  \to
  (u_{T-1},X_T,m_{T-1}),
\]
so gradients flow through both the PDE state and the self-talk state. At
the next optimizer update, the numeric values of \(X_T\) and \(m_{T-1}\)
are carried forward, but the graph is detached:
\[
  X_T\leftarrow \operatorname{stopgrad}(X_T),
  \qquad
  m_{T-1}\leftarrow \operatorname{stopgrad}(m_{T-1}).
\]
This preserves the one-way ecological time stream while keeping each
optimization problem finite and memory-bounded.

\section{Differentiable PDE Tooling}

The first implementation should use a PyTorch differentiable age-grid
solver. Let
\[
  a_j=j\Delta a,\qquad j=0,\ldots,N_a.
\]
Use a first-order upwind update for the transport term:
\begin{equation}
  \frac{x_{i,j}^{n+1}-x_{i,j}^{n}}{\Delta t}
  +
  \frac{x_{i,j}^{n}-x_{i,j-1}^{n}}{\Delta a}
  =
  -x_{i,j}^{n}\,\mathcal R_{i,j}^{n},
  \qquad j\ge 1,
  \label{eq:upwind}
\end{equation}
where \(\mathcal R_{i,j}^{n}\) contains mortality, control, and
interaction terms. The boundary is updated by quadrature:
\begin{equation}
  x_{i,0}^{n+1}
  =
  \sum_{j=0}^{N_a}
  k_i(a_j,t_n)x_{i,j}^{n}\Delta a.
  \label{eq:boundary}
\end{equation}

All state updates, quadratures, \(\eta\) computation, softplus control,
and losses should remain in PyTorch tensors so that the equilibrium loss
can backpropagate through the rollout to the controller.

\section{Time-Varying Ecology}

The ecology should not be static. The first version uses the same profile
families used in the Lotka--Sharpe neural-operator experiments
\cite{krstic2026lotkasharpe}:
\begin{align}
  k(a,t)
  &=
  k_{\rm base}
  +
  k_{\rm amp}(t)
  \exp\left(
    -\frac{(a-k_{\rm center})^2}{2k_\sigma^2}
  \right),
  \\
  \mu(a,t)
  &=
  \mu_{\rm base}(t)
  +
  \mu_{\rm juv,amp}e^{-\mu_{\rm juv}a}
  +
  \mu_{\rm sen,amp}(t)a^{\mu_{\rm sen}},
  \\
  g(a)
  &=
  g_{\rm base}
  +
  g_{\rm amp}
  \exp\left(
    -\frac{(a-g_{\rm center})^2}{2g_\sigma^2}
  \right).
\end{align}
The centers, widths, and interaction profile are fixed in the first
version. The ecological regime changes through mixed-sin modulation of
amplitude-like parameters:
\[
  k_{\rm amp}(t),\qquad
  \mu_{\rm base}(t),\qquad
  \mu_{\rm sen,amp}(t).
\]

A generic mixed-sin modulation is
\begin{equation}
  q(t)
  =
  q_0
  \left[
    1+
    \rho
    \frac{\sum_{\ell=1}^{L} A_\ell\sin(\omega_\ell t+\varphi_\ell)}
         {\sum_{\ell=1}^{L}|A_\ell|}
  \right],
  \qquad 0<\rho<1.
  \label{eq:mixed-sin}
\end{equation}
The model does not observe \(q(t)\), \(k(a,t)\), \(\mu(a,t)\), or
\(\zeta(t)\). It only observes \(\eta_1(t)\), \(\eta_2(t)\), and the
constant pulse. The language loop is therefore tested as an internal
state that may help track hidden ecological regime changes.

\section{Continuous Training and Testing}

Training should follow the exp0522 continuous case:
\begin{itemize}[leftmargin=2em]
  \item time moves forward in one long ecological stream;
  \item training windows are consecutive pieces of this stream;
  \item the PDE state, language state, and error view are carried forward;
  \item the carried states are detached at window boundaries;
  \item the stream is not rewound to \(t=0\) after each update.
\end{itemize}

Evaluation should reserve a future segment of the same ecological stream.
The model trains on the earlier segment, then rolls into the held-out
future segment without weight updates. This tests whether the learned
closed-loop controller continues to stabilize the ecology as the hidden
birth and mortality parameters keep evolving.

\section{First Experiment Design}

The first experiment should compare:
\begin{enumerate}[leftmargin=2em]
  \item \textbf{full self-talk controller:}
  \(\eta\), pulse, error view, and language all live;
  \item \textbf{no-language controller:}
  same external observations, but \(m_t\equiv 0\);
  \item \textbf{no-error-view controller:}
  language live, but \(e_t\equiv 0\);
  \item \textbf{minimal controller:}
  no language and no error view, only \((\eta_1,\eta_2,1)\);
  \item \textbf{blink/stutter rollouts:}
  temporary removal of error or language during evaluation.
\end{enumerate}

The primary metrics are:
\begin{itemize}[leftmargin=2em]
  \item mean and final \(\eta_1^2+\eta_2^2\);
  \item mean control energy \(u^2\);
  \item positivity of \(x_i(a,t)\) and \(u(t)\);
  \item future-segment stability after training stops;
  \item recovery after blink/stutter interruptions.
\end{itemize}

\section{Working Summary}

The paper supplies a biologically meaningful full PDE environment and
equilibrium coordinates. Exp0522 supplies the self-talk MLP and continuous
rollout protocol. Exp0526 combines them by training a controller directly
from closed-loop equilibrium loss:
\[
  \big(\eta_1(t),\eta_2(t),1;\text{self-talk}\big)
  \mapsto
  \tilde u_t
  \mapsto
  u_t=\softplus(\tilde u_t)
  \mapsto
  x(\cdot,t+1)
  \mapsto
  \eta(t+1)
  \mapsto
  \mathcal L_\eta.
\]

\bibliographystyle{plain}
\bibliography{references}

\end{document}
